\chapter{README}

章タイトルのデザインは quotchap パッケージ\footnote{\url{https://ctan.org/pkg/quotchap}}のパクリ。

\section{執筆のしかた}

このリポジトリを \verb`git clone` したら、\verb`*.tex` をてきとーに編集して、\verb`make` するだけ。

必要に応じて、\verb`git remote remove origin` してから自前のリポジトリに \verb`git remote add origin ...` しなおすべし。

\subsection{\LaTeX 環境}

Lua\LaTeX + jlreq が必要。\TeX Live で動くはず。

docker で構築するのが楽ちん。

\begin{lstlisting}[style=tty]
% docker pull texlive/texlive:latest
\end{lstlisting}

docker を使わない場合は \verb`Makefile` の \verb`LATEXMK=...` の行を修正する必要がある。

\begin{lstlisting}
# docker を使う場合
LATEXMK= docker run --rm -v .:/workdir texlive/texlive:latest latexmk
# 使わない場合
#LATEXMK= latexmk
# ローカルではdockerを使わないが、github actions でコンパイルするときには使う
#LATEXMK= $(shell command -v latexmk || echo "docker run --rm -v .:/workdir texlive/texlive:latest latexmk")
\end{lstlisting}

\subsection{コンパイル}

\verb`make` 一発。

\verb`Makefile` をとくに修正していない場合、電子版(画面上での閲覧を想定したPDF)の \verb`main.pdf` を生成する。また、\verb`make all` ではそれに加えて印刷版(印刷所への入稿を想定したPDF)の \verb`print.pdf` を生成する。コンパイルのたびに両方生成するのは時間のムダなので、デフォルトではあえて片方だけ生成するようにしている。

2種類のPDFを生成するが、電子版、印刷版どちらに対する変更も \verb`main.tex` に対しておこなう。\verb`print.tex` は電子版か印刷版かを区別するためのマクロを定義してから \verb`main.tex` を読み込むだけのファイルで、通常はとくに修正する必要はない。

以下のように記述すると電子版、印刷版で処理を切り替えられる。

\begin{lstlisting}
\ifdefined\PRINT
  % 印刷版の処理
\else
  % 電子版の処理
\fi
\end{lstlisting}

\verb`make` の際、\verb`git` のコミットハッシュが記載された \verb`revid.tex` というファイルを自動で生成する。この値は \verb`\revid` というマクロで参照できる。本ドキュメントでは奥付で使用している。また、PDFのメタ情報にも埋め込んでいる。

\begin{lstlisting}[style=tty]
% git rev-parse --short main
1eb445e
% cat revid.tex
\def\revid{1eb445e}
% pdfinfo -meta main.pdf | fgrep VersionID
                  <pdfaProperty:name>VersionID</pdfaProperty:name>
      <xmpMM:VersionID>1eb445e</xmpMM:VersionID>
\end{lstlisting}

\subsection{github workflow}

\verb`main` または \verb`master` ブランチで github に push すると、github 上で \verb`make all` される。docker を使わずにコンパイルするように \verb`Makefile` の \verb`LATEXMK=...` の行を修正していた場合、github でのビルドがコケるので注意。

自動コンパイルしたくない場合は \verb`.github/` 以下をまるごと削除すべし。

\section{セクション}

ここからは見た目を確認するためのデモ。

\subsection{サブセクション}

びよーん。

\subsubsection{サブサブセクション}

ぱおーん。

\begin{column}{コラムタイトル}

コラムはこう書く\footnote{コラム内で脚注を書くとこうなる}。

目次にも自動的に掲載される。

\end{column}

ソースコードはリスト\ref{src}のように書く。

\begin{lstlisting}[caption=hello.sh, label=src]
echo Hello, world!
\end{lstlisting}

実行するとこうなる。

\begin{lstlisting}[style=tty]
$ sh hello.sh
Hello, world!
\end{lstlisting}


以下は \verb`framed` パッケージで提供されていた環境を \verb`tcolorbox` でそれっぽく作ったもの。

なんか引用文を入れるのによさそうな環境。えらい人曰く、

\begin{leftbar}
腹へった。
\end{leftbar}

用途不明だけどなんか枠。

\begin{snugshade}
グレーの箱。
\end{snugshade}

